\documentclass{article}

% NeurIPS 2026 style
\usepackage[preprint]{neurips_2025}
\usepackage[utf8]{inputenc}
\usepackage[T1]{fontenc}
\usepackage{hyperref}
\usepackage{url}
\usepackage{booktabs}
\usepackage{amsfonts}
\usepackage{amsmath}
\usepackage{microtype}
\usepackage{graphicx}
\usepackage{xcolor}
\usepackage{multirow}
\usepackage{subcaption}

\title{CIPHER: A Benchmark for Captcha Instruction Parsing with Hidden Extraction and Reasoning}

\author{%
  Maxim Kurkin \\
  \texttt{kurkin@2a2i.org}
}

\begin{document}
\maketitle

% ============================================================
\begin{abstract}
We introduce \textbf{CIPHER} (\textbf{C}aptcha \textbf{I}nstruction \textbf{P}arsing with
\textbf{H}idden \textbf{E}xtraction and \textbf{R}easoning), a procedurally generated benchmark
for evaluating vision-language models (VLMs) on video-editing captcha tasks.
In CIPHER, an agent receives a short video clip together with natural-language editing
instructions (e.g.\ ``reverse the video, then flip horizontally'') and must
(i)~parse the instructions into a sequence of edit operations,
(ii)~execute those operations via \texttt{ffmpeg}, and
(iii)~extract a short alphanumeric key that becomes visible in the correctly edited result.
The benchmark spans five difficulty levels that independently vary key visibility,
operation count, and instruction ambiguity, yielding \num{250} procedurally generated samples
backed by diverse Kinetics-400 background clips.
We evaluate OCR-only baselines and a suite of frontier VLMs, finding that
\ldots % TODO: fill in main finding
CIPHER reveals a clear capability gap between stages: instruction parsing saturates early
while key extraction from realistic backgrounds remains an open challenge.
\end{abstract}

% ============================================================
\section{Introduction}
\label{sec:intro}

Captchas have evolved from distorted text images to dynamic, multi-step challenges
designed to resist automation while remaining trivially solvable by humans.
A recent class of video-editing captchas presents users with a short clip and a natural-language
instruction such as ``reverse the video, then flip it horizontally.''
Following these instructions reveals an embedded alphanumeric key that serves as the captcha answer.
A human solves this in seconds; an automated agent must parse the instruction,
execute the correct sequence of video transformations, and extract the now-visible key---
a compound task spanning language understanding, procedural execution, and visual perception.

Despite the proliferation of video-capable VLMs~\citep{gpt4o,gemini,qwen2vl},
no benchmark systematically tests this compound capability.
Existing video benchmarks target temporal reasoning~\citep{videomme,longvideobench}
or egocentric understanding~\citep{ego4d}, not instruction-conditioned video manipulation.
Video-editing benchmarks~\citep{vbench,editbench} evaluate generative models rather than agents.
CIPHER fills this gap.

\paragraph{Contributions.}
\begin{enumerate}
  \item \textbf{Task formulation.} We define the video-editing captcha solving task
        as a three-stage pipeline: instruction parsing (S1), ffmpeg execution (S2),
        and key extraction (S3), enabling fine-grained per-stage diagnostics.

  \item \textbf{Benchmark.} CIPHER provides 250 procedurally generated samples
        across five difficulty levels (L1--L5), varying key visibility,
        operation count (1--4 ops), and instruction ambiguity (exact to compositional).
        Backgrounds are drawn from Kinetics-400 to ensure visual diversity.

  \item \textbf{Metrics.} We propose \emph{Pipeline-F1}, a weighted combination of
        per-stage scores ($0.25 \cdot \text{S1} + 0.25 \cdot \text{S2} + 0.50 \cdot \text{S3}$),
        alongside strict exact match (EM) and character-level F1.

  \item \textbf{Baselines.} We benchmark OCR-only, Claude Haiku, GPT-4o-mini,
        Gemini Flash, and local open-weight models (Qwen2-VL-7B/72B),
        revealing that S1 parsing saturates quickly while S3 extraction
        on natural backgrounds remains the binding bottleneck.
\end{enumerate}

\section{Task Definition}
\label{sec:task}

\paragraph{Setting.}
Let $\mathcal{V}$ be the space of video clips and $\mathcal{L}$ the space of
natural-language strings. CIPHER evaluates a function
$f: \mathcal{V} \times \mathcal{L} \to \Sigma^7$ that maps a corrupted video
$v'$ and an instruction $\mathcal{I}$ to a predicted target string $\hat{k}$.
The original clip $v$ contains a latent key $k^* \in \Sigma^7$ embedded during
a temporal window $[t_s, t_e]$. A sequence of transformation operations
$\mathbf{o} = (o_1, \dots, o_n)$ is applied to $v$ to yield
$v' = o_n \circ \dots \circ o_1 (v)$. The agent's objective is to recover the
inverse program $\mathbf{o}^{-1}$ from $\mathcal{I}$ and $v'$, execute it, and
extract $k^*$ from the recovered terminal state.

The key $k^*$ is drawn from unambiguous alphanumerics
($\Sigma = \text{A--Z} \cup \text{2--9} \setminus \{\text{O,0,I,1}\}$).
Instructions are generated from $\mathbf{o}$ via paraphrase templates.

\paragraph{Three-stage pipeline.}
This task is intentionally composite. An end-to-end failure
($\hat{k} \neq k^*$) may arise because the model inferred the wrong program,
realized the wrong tool chain, or failed to localize and read the key after
successful recovery. By enforcing a three-stage pipeline---Instruction Parsing
(S1), Execution (S2), and Key Extraction (S3)---CIPHER isolates these failure
modes. In policy terms, S1 corresponds to \emph{program inference}, S2 to
\emph{tool-chain realization}, and S3 to \emph{terminal-state extraction}.
This decomposition prevents monolithic metrics from masking which part of the
reasoning policy is actually failing.

\begin{description}
  \item[S1 Parse:] $\mathcal{I} \to \hat{\mathbf{o}}$. Score: sequence-aware op F1.
  \item[S2 Execute:] run \texttt{ffmpeg}($\hat{\mathbf{o}}$). Score: binary validity.
  \item[S3 Extract:] read $\hat{k}$ from edited video. Score: CharF1 + EM.
\end{description}

\paragraph{Operations.}
The operation vocabulary is intentionally minimal: six composable
\texttt{ffmpeg} transforms, \texttt{reverse}, \texttt{hflip},
\texttt{vflip}, \texttt{rotate90}, \texttt{rotate270}, and
\texttt{speed2x}. These are reversible, visually interpretable, and cover both
spatial and temporal recovery. The restricted vocabulary is a design choice:
it keeps the benchmark focused on inverse program recovery rather than on open
ended action-space exploration.

\paragraph{Video-temporal diagnostic domain.}
In addition to the original inverse-\texttt{ffmpeg} pipeline, we instantiate a
controlled video-temporal split-key domain for faster iteration on temporal
aggregation. A key is split into two fragments that appear on two consecutive
frames of a short clip sampled from the public DAVIS-2017 video corpus
\citep{davis2017}. The solver must identify the target frames, read both
fragments, and invert a small symbolic operation before answering. The
operation set is \texttt{none}, \texttt{swap\_fragments},
\texttt{reverse\_chars}, and \texttt{swap\_reverse\_chars}; the last two
require character-level transformation rather than simple OCR.

This domain deliberately separates \emph{temporal evidence acquisition} from
generic video editing: the latent program, target frame indices, fragment
bounding boxes, displayed fragments, normalized fragments, and final answer are
all known from generation. This enables dense process scoring. Besides strict
EM, we assign partial credit for operation recognition, target-frame
localization, displayed-fragment similarity, normalized-answer similarity,
transform-derived answer similarity, final-answer similarity, and successful
tool use. The resulting traces are used both for model diagnosis and for
evolutionary seed selection.

\section{The CIPHER Benchmark}
\label{sec:benchmark}

\subsection{Difficulty levels}

\begin{table}[h]
\centering\small
\caption{CIPHER difficulty levels (50 samples each, $N{=}250$ total).}
\label{tab:levels}
\begin{tabular}{lcccp{5cm}}
\toprule
Level & Key vis. & Ops & NL ambiguity & Primary challenge \\
\midrule
L1 & High   & 1 & Exact         & Single-op baseline \\
L2 & Medium & 2 & Exact         & Two-op chaining \\
L3 & Low    & 2 & Paraphrase    & Surface-form variation \\
L4 & Low    & 3 & Reversed      & Reversed description \\
L5 & Min.   & 4 & Compositional & Underspecified instruction \\
\bottomrule
\end{tabular}
\end{table}

\paragraph{Key visibility} varies font size (12\%$\to$5\% of frame height),
alpha (255$\to$140), and exposure duration (4.5$\to$2.5s).

\paragraph{NL ambiguity templates:}
(1)~exact canonical; (2)~paraphrase from pool of 5 per op;
(3)~reversed order with ``undo'' framing;
(4)~compositional: ``The video has been edited. Undo all transformations.''

\paragraph{Operation Vocabulary.}
The operational vocabulary consists of seven foundational transformations: \texttt{reverse}, \texttt{hflip}, \texttt{vflip}, \texttt{rotate90}, \texttt{rotate270}, \texttt{speed2x}, and \texttt{grayscale}. This specific set was chosen because these operations are visually distinct, non-destructive (except \texttt{grayscale}, used primarily as a distractor), and strictly composable via standard \texttt{ffmpeg} filters. They demand different types of reasoning: spatial (\texttt{hflip}, \texttt{rotate}), temporal (\texttt{reverse}, \texttt{speed2x}), and color-based (\texttt{grayscale}), providing a comprehensive testbed for multimodal execution without requiring complex parameter tuning from the model.

\subsection{Procedural generation and Dataset Integrity}

Samples are fully deterministic from $(seed, level)$.
Background clips are drawn uniformly from a Kinetics-400 pool
(diverse human activities, minimal on-screen text).
Test split: seeds 0--49. Challenge split: seeds 50--99.

A persistent challenge in evaluating large language and vision models is the risk of data contamination, where static benchmark data inadvertently leaks into the pre-training corpus. CIPHER mitigates this risk entirely through programmatic dataset synthesis. The benchmark is procedurally generated from a tuple of $(seed, \text{level})$, which deterministically selects the background clip, the randomly generated 7-character key, the sequence of operations, and the natural language instruction template. Because samples are dynamically constructed on-demand using disjoint background video pools, researchers can regenerate the benchmark with novel seeds, guaranteeing an uncontaminated evaluation environment and an effectively infinite supply of test instances.

\subsection{Metrics}
\label{sec:metrics}

\paragraph{EM.} $\frac{1}{N}\sum_i \mathbf{1}[\hat{k}_i = k^*_i]$.

\paragraph{CharF1.} Character-level precision/recall F1 (partial credit).

\paragraph{PipeF1.}
\begin{equation}
\text{PipeF1} = 0.25\cdot\overline{\text{S1}} + 0.25\cdot\overline{\text{S2}} + 0.50\cdot\overline{\text{CharF1}}
\end{equation}
S3 receives double weight as the terminal objective;
S1/S2 serve as diagnostic leading indicators.
\section{Experiments}
\label{sec:experiments}

\subsection{Experimental setup}

\paragraph{Inference modes.}
We evaluate models in two settings:
(1)~\textbf{S3-only}: models receive a ground-truth corrected frame containing
the key and must read it. This tests pure terminal-state extraction and serves
as an upper-bound control slice.
(2)~\textbf{Full pipeline}: models receive the scrambled \texttt{.mp4} video
and natural-language instructions with Bash access, and must independently
parse instructions (S1), execute \texttt{ffmpeg} corrections (S2), locate the
key window, and extract the key (S3).

The distinction is central: strong performance in S3-only does not imply strong
inverse program recovery. It only shows that once the correct terminal state is
provided, the model can decode the target string. Full-pipeline performance
measures whether the system can actually reach that terminal state.

\paragraph{Agentic evaluation.}
For full-pipeline evaluation, VLM agents are given Bash access and must decide
their own video processing strategy (ffmpeg flags, frame sampling, temporal
search). No ground-truth metadata (operations, key, timing) is provided.

\paragraph{Reported metrics.}
We report strict exact match (EM), character-level F1, stage-specific
S1 parse F1, S2 execution success rate, and Pipeline-F1. In the current draft,
the strongest supported evidence comes from S3-only and calibrated extraction
experiments; full-pipeline reporting remains incomplete and is therefore
presented as an explicit gap rather than as a filled leaderboard.

For the video-temporal diagnostic domain, we additionally report \emph{process
partial credit}. This score decomposes a failed answer into operation credit,
frame-localization credit, displayed-fragment similarity, normalized-answer
similarity, transform-derived-answer similarity, final-answer similarity, and
tool-use credit. It is intentionally not a replacement for EM; it is a dense
diagnostic signal used to distinguish uninformative failures from meaningful
near misses.

\subsection{Models}

We benchmark OCR baselines, heuristic pipelines, and frontier multimodal
systems:

\begin{description}
  \item[\textbf{OCR-only}]
  Pure extraction baseline. Uses easyocr on sampled key frames; no language
  understanding, program inference, or tool use.

  \item[\textbf{Yellow-OCR}]
  Color-filtered OCR baseline. Isolates yellow pixels via HSV masking before
  running easyocr. Tests whether color-aware preprocessing can substitute for
  reasoning or executable recovery.

  \item[\textbf{MCP Skill}]
  Heuristic agentic baseline (\texttt{solve\_captcha.py}): regex-based
  instruction parsing (S1), \texttt{ffmpeg} execution (S2), and a multi-method
  extraction chain (S3).

  \item[\textbf{Claude Haiku / Sonnet / Opus}]
  API-first frontier multimodal systems used to probe the gap between prepared
  extraction and end-to-end recovery.

  \item[\textbf{GPT-5.4}]
  Tool-using frontier model evaluated through the same full-pipeline protocol.
\end{description}

\subsection{Current results}

\begin{table}[h]
\centering\small
\caption{S3 key extraction from ground-truth corrected frames ($n{=}5$ per
level). Best per column in \textbf{bold}.}
\label{tab:s3}
\begin{tabular}{lrrrr|r}
\toprule
Model & L1 & L2 & L3 & L4 & Total \\
\midrule
OCR-only       & 20\% & 25\% & 30\% & 25\% & 25\% \\
Yellow-OCR     & 33\% & --   & --   & --   & --   \\
MCP Skill      & 20\% & 25\% & 30\% & 25\% & 25\% \\
Haiku          & \textbf{100\%} & 60\% & 40\% & \textbf{100\%} & 75\% \\
Sonnet         & \textbf{100\%} & \textbf{100\%} & \textbf{100\%} & \textbf{100\%} & \textbf{100\%} \\
Opus           & \textbf{100\%} & \textbf{100\%} & \textbf{100\%} & \textbf{100\%} & \textbf{100\%} \\
\bottomrule
\end{tabular}
\end{table}

\begin{table}[h]
\centering\small
\caption{Full pipeline results (video $\to$ key). This table is intentionally
left incomplete in the current draft; filling it is a required next experiment
because the benchmark's main claim concerns the gap between S3-only extraction
and end-to-end recovery.}
\label{tab:main}
\begin{tabular}{lrrrrrr}
\toprule
Model & L1 EM & L2 EM & L3 EM & L4 EM & EM (all) & PipeF1 \\
\midrule
MCP Skill     & -- & -- & -- & -- & -- & -- \\
Haiku         & -- & -- & -- & -- & -- & -- \\
Sonnet        & -- & -- & -- & -- & -- & -- \\
Opus          & -- & -- & -- & -- & -- & -- \\
GPT-5.4       & -- & -- & -- & -- & -- & -- \\
\bottomrule
\end{tabular}
\end{table}

\paragraph{Overview.}
Table~\ref{tab:s3} reports S3 extraction accuracy in isolation: each model
receives a corrected frame containing the yellow key and must read it. Frontier
VLMs (Opus, Sonnet) achieve perfect accuracy, while Haiku drops on smaller,
randomly positioned keys. OCR baselines perform poorly due to background text
interference and clutter.

Crucially, S3 in isolation is \emph{not} the binding bottleneck for frontier
models. The benchmark's real difficulty lies in the full recovery policy:
agents must infer the inverse transformation program, execute it correctly, and
identify the correct temporal window containing the key. The missing full
pipeline table is therefore not a cosmetic omission; it is the central
remaining experiment required to validate the benchmark's main claim.

\subsection{Video-temporal diagnostic results}

We also ran a controlled DAVIS-2017 video-temporal split in which the model
must aggregate two consecutive frames and invert a symbolic fragment
transformation. This split is not yet the full inverse-\texttt{ffmpeg}
benchmark; it is the current fast iteration domain for studying temporal
aggregation, solvability, and dense process metrics. The current pool contains
78 generated tasks with process metrics and 233 model responses across three
CLI backends. All modern samples in this slice use dynamic DAVIS frame
sequences and pass the current ``likely-readable'' solvability gate.

\begin{table}[h]
\centering\small
\caption{Current DAVIS video-temporal diagnostic slice. EM is strict exact
match. Missing counts empty/NONE final answers. Partial is dense process
partial credit. Values are means over response rows; $\pm$ is population
standard deviation.}
\label{tab:video-temporal-current}
\begin{tabular}{lrrrr}
\toprule
Backend & $n$ & EM $\uparrow$ & Missing $\downarrow$ & Partial $\uparrow$ \\
\midrule
Claude CLI Haiku & 78 & $14.1{\pm}34.8$ & $1.3{\pm}11.2$ & $43.1{\pm}28.6$ \\
Codex CLI GPT-5.3 Spark & 77 & $0.0{\pm}0.0$ & $100.0{\pm}0.0$ & $6.5{\pm}5.6$ \\
Gemini CLI Flash Lite & 78 & $0.0{\pm}0.0$ & $98.7{\pm}11.2$ & $5.0{\pm}5.7$ \\
\bottomrule
\end{tabular}
\end{table}

At the task level, 11 of 78 tasks are \emph{strict-good}: at least one model
achieves EM. A larger set of 45 tasks is \emph{semantic-good}, meaning either
EM or maximum process partial credit at least 0.25. Finally, 77 tasks have at
least one non-missing response with partial credit at least 0.15, making them
usable for debugging even when no backend solves them exactly. The gap between
strict-good and semantic-good is the most important observation: the split
contains many failures that are not random, but partially correct trajectories.

The strongest current backend is Claude CLI Haiku, which often identifies the
operation and produces structured intermediate evidence but still loses
characters or over-reads distractors on hard samples. Codex Spark and Gemini
Flash Lite mostly fail by returning no final answer in this setup, although
their traces occasionally receive small operation/tool credit. This suggests
that backend/tool framing and response protocol are part of the benchmark
surface, not merely implementation details.

\subsection{Required experiments for the camera-ready version}

The current draft still lacks the analyses that best express CIPHER's identity
as a diagnostic benchmark:
\begin{itemize}
  \item scaling curves over chain length, instruction ambiguity, distractor
        density, and temporal sparsity;
  \item paired S3-only vs full-pipeline comparisons with an explicit
        execution-gap metric;
  \item stage-attribution plots for S1, S2, and S3 across model families; and
  \item a curated failure-seed suite with traces.
\end{itemize}

The video-temporal slice supplies the first version of this failure-seed suite:
12 seedbank genomes were selected by validity, non-missing rate, process
partial credit, frame-localization credit, transform similarity, final-answer
similarity, and diversity. The selected genomes cover all four fragment
operations and range from easy validity-tuned settings (4--6 frames, low noise,
large bright text) to hard-mined settings (7--8 frames, high jitter, distractors,
and \texttt{swap\_reverse\_chars}).

\section{Analysis}
\label{sec:analysis}

CIPHER is most valuable when failures are treated as structured signals rather
than as a flat collection of incorrect strings. The analysis should therefore
be read as a \emph{failure geometry}: which stage breaks first, how errors
change with difficulty, and which seeds should be preserved as regression
cases. Some subsections below remain partially prospective pending
full-pipeline results, but the intended analytical frame is already clear.

\subsection{Stage bottleneck analysis}

We decompose failures by stage to identify where the recovery policy breaks
down.

\paragraph{S1 (program inference).}
On simple levels, S1 often saturates first: regex matching achieves high parse
F1 on canonical instructions. The remaining errors are not random; they
concentrate in paraphrastic or compositional instructions where the system must
recover the latent operator sequence rather than match surface forms.

\paragraph{S2 (tool-chain execution).}
Execution failures are rare on short chains but increase with operation count,
primarily because the model must realize a valid \texttt{ffmpeg} chain rather
than merely name the right operations.

\paragraph{S3 (terminal-state extraction).}
S3 is \emph{not} the binding bottleneck for frontier VLMs once the correct
frame is provided. The true difficulty lies in the coupling between temporal
localization and execution: many systems that can read the key still fail to
arrive at the correct frame.

\paragraph{Near-miss errors are informative.}
Near-miss predictions should not be discarded as noise. They reveal whether a
model reached the correct terminal state but misread one character, localized
the wrong frame, or recovered the wrong but visually plausible distractor. In a
diagnostic benchmark, these distinctions matter more than a single binary EM
label.

\subsection{Character-level confusion in the near-miss regime}

Under calibrated extraction settings, models produce consistent substitutions
among visually similar glyphs rather than random strings. These errors are
useful because they identify a readable-but-fragile regime in which the task is
still discriminative.

\subsection{Video-temporal failure taxonomy}

The DAVIS video-temporal slice makes the failure geometry concrete. Among the
78 modern video-temporal tasks, 67 currently have no exact match from any
evaluated backend, but only one task is an all-missing failure across all
models. This is the desired hard-but-diagnostic regime: most failures still
contain usable intermediate evidence.

The dominant failure families are:
\begin{itemize}
  \item \textbf{near misses}: 21 tasks have no EM but reach maximum process
        partial credit at least 0.5, indicating that a model recovered part of
        the operation/frame/fragment structure but failed the final string;
  \item \textbf{frame-localization weakness}: 24 tasks have maximum
        frame-localization credit below 0.5, showing that models often read
        from non-target frames or from the entire clip rather than the two
        target frames;
  \item \textbf{over-reading}: 12 tasks contain predictions longer than the
        gold key by more than two characters, usually because the model absorbs
        distractors or neighboring yellow text into the answer; and
  \item \textbf{hard swap-and-reverse cases}: the highest-difficulty frontier
        is dominated by \texttt{swap\_reverse\_chars} seeds with 7--8 frames,
        high jitter, distractors, and non-central placement.
\end{itemize}

This taxonomy directly motivates the dense metric. A flat EM score would mark
all of these examples as identical failures. Process partial credit separates
``no evidence'', ``found the operation'', ``localized the wrong frame'', and
``read the right evidence but transformed it incorrectly'', which are different
failure modes for an agentic multimodal system.

\subsection{Canonical failure suite}

For benchmark-driven development, we recommend preserving a curated suite of
20--50 canonical failure seeds with associated traces. Each entry should store:
seed, level, background regime, gold operator chain, predicted operator chain,
execution trace, temporal localization decision, predicted key, and error
family. This turns CIPHER from a static benchmark into a replayable regression
environment for agentic systems.

The current video-temporal seedbank is the first implementation of this idea.
It stores not only the winning candidates, but also diverse non-winning
candidates selected by a utilization--diversity criterion. This matters because
the most useful benchmark samples are not always the hardest candidates by EM:
some are slightly easier but much more interpretable, and some expose specific
operation or temporal-localization weaknesses that would be lost if we kept
only the scalar best-of-run candidate.

\section{Conclusion}
\label{sec:conclusion}

We introduced CIPHER, a procedurally generated benchmark for video-editing captcha solving
that stress-tests VLMs on a compound task: parsing natural-language edit instructions,
executing an ffmpeg transformation chain, and extracting a hidden key from the result.
The three-stage pipeline structure and PipeF1 metric enable fine-grained diagnosis
of where agent capability breaks down.

Our experiments reveal that frontier VLMs solve S3 key extraction near-perfectly
in isolation (Opus/Sonnet: 100\% EM), while OCR baselines struggle ($\sim$25\%).
The compound difficulty of CIPHER lies in the full agentic pipeline:
temporal key localization, correct operation sequencing, and tool execution
create a challenging integration problem that no evaluated model solves reliably
across all levels.

\paragraph{Limitations.}
CIPHER focuses on seven deterministic ffmpeg operations;
generative edits (e.g., inpainting, style transfer) and audio-only keys are out of scope.
Human baselines are not included in this version.

\paragraph{Future work.}
L6 (QR/morse key encoding), multi-agent solving, and adversarial instruction generation
are natural extensions. The procedural generator enables zero-cost dataset refresh,
making CIPHER suitable for ongoing leaderboard tracking.

\paragraph{Reproducibility.}
All code and generation scripts are released at
\url{https://github.com/Fr0do/cipher}.
The test manifest (seeds 0--49) is pinned at commit \texttt{TODO}.


\section*{Acknowledgements}

\bibliography{references}
\bibliographystyle{plainnat}

\end{document}
